Queda por demostrar (6.14.1.3). Es evidente que
$A'\subset K'\cap\overline{A}'_{\mathfrak p'}$. Recíprocamente, sea
$y'$ un elemento de esta intersección. Sea $a$ un «parámetro» para
$A$, de modo que $Aa$ sea $\mathfrak m$-primario, y sea $a'$ la imagen
de $a$ en $A'$. Todo elemento de $K$ puede escribirse $x/a^n$ con
$x\in A$ y $n>0$, pues $Aa$ contiene una potencia de $\mathfrak m$;
por tanto, se puede escribir $y'=x'/a'^n$ con $x'\in A'$. Observemos
ahora que $\mathfrak p'$ es el \emph{único} ideal primo asociado a
$A'/a'^nA'=(A/a^nA)\otimes_A A'$. Como $A'$ es un $A$-módulo plano,
esto se deduce de (3.3.1): $\mathfrak m$ es el único ideal primo
asociado a $A/a^nA$, y $k\otimes_A A'$ es un dominio entero. En
consecuencia, $a'^nA'$ es la imagen inversa en $A'$ de
$a'^nA'_{\mathfrak p'}$. Como, por hipótesis,
$y'\in\overline{A}'_{\mathfrak p'}$, la imagen de $x'$ en
$A'_{\mathfrak p'}$ pertenece a $a'^nA'_{\mathfrak p'}$; de ahí
$x'\in a'^nA'$ y $y'\in A'$, lo que termina la demostración.
\end{proof}

